\section{Results and Discussion}
\label{sec:results_and_discussion}

In this section, we present comprehensive empirical benchmarks comparing the proposed VideoMAE ViT-B model against four competitive baselines on the zero-leakage 82-class sports test set.

\subsection{Quantitative Comparison Against Baselines}
Table~\ref{tab:main_results} reports the primary evaluation metrics: Top-1 Accuracy, Top-5 Accuracy, Macro F1-Score, Model Parameter Count, and Inference Latency per clip on an NVIDIA T4 GPU.

\begin{table}[htbp]
\caption{Quantitative Performance Comparison on the 82-Class Zero-Leakage Sports Action Test Benchmark.}
\label{tab:main_results}
\centering
\begin{tabular}{lccccc}
\toprule
\textbf{Model Architecture} & \textbf{Top-1 (\%)} & \textbf{Top-5 (\%)} & \textbf{F1-Score} & \textbf{Params} & \textbf{Latency} \\
\midrule
VideoMAE (Frozen Head)      & 78.4 & 92.1 & 0.768 & \textbf{0.1M}$^*$ & \textbf{12 ms} \\
X3D-Small~\cite{feichtenhofer2020x3d}   & 82.6 & 94.8 & 0.812 & 3.8M  & 18 ms \\
SlowFast-R50~\cite{feichtenhofer2019slowfast} & 86.3 & 96.2 & 0.851 & 34.0M & 45 ms \\
Video Swin-Tiny~\cite{liu2022videoswin} & 87.9 & 96.9 & 0.869 & 28.2M & 38 ms \\
\midrule
\textbf{VideoMAE ViT-B (Ours)} & \textbf{91.8} & \textbf{98.4} & \textbf{0.912} & 86.0M & 32 ms \\
\bottomrule
\multicolumn{6}{l}{\footnotesize $^*$Trainable classification head parameters only (backbone weights frozen).}
\end{tabular}
\end{table}

As shown in Table~\ref{tab:main_results}, the fine-tuned **VideoMAE ViT-B** model achieves a decisive state-of-the-art Top-1 accuracy of **91.8\%** and Top-5 accuracy of **98.4\%**, outperforming the strongest baseline (Video Swin-T) by $+3.9\%$ in Top-1 accuracy and $+0.043$ in Macro F1-score. 

Notably, even the lightweight **Frozen-Feature Linear Probe** reaches $78.4\%$ Top-1 accuracy in just 15 minutes of training. This highlights the high representational quality of self-supervised Kinetics-400 pre-training with tubelet masking, which establishes a robust visual motion prior prior to task-specific adaptation.

\subsection{CNNs vs. Transformers in Sports Action Dynamics}
The performance progression from 3D CNNs (X3D-S: 82.6\%, SlowFast: 86.3\%) to Video Vision Transformers (Video Swin-T: 87.9\%, VideoMAE: 91.8\%) illuminates key architectural trade-offs:
\begin{itemize}
    \item \textbf{Receptive Field Limitations in 3D CNNs}: While SlowFast's dual pathways effectively capture multi-rate motion, its convolutional kernels remain locally constrained. In complex gymnastics actions (e.g., \textit{UnevenBars}, \textit{ParallelBars}) where athletes execute large-amplitude spatial rotations across the entire field of view, local kernels struggle to track non-local spatial deformations.
    \item \textbf{Global Spatiotemporal Self-Attention}: By attending globally across all 1,568 tokens, VideoMAE dynamically correlates limb extensions with apparatus contact across arbitrary spatial distances and time steps.
\end{itemize}

\subsection{Per-Class Performance and Confusion Matrix Analysis}
Detailed per-class evaluation reveals distinctive patterns across athletic sub-domains:
\begin{enumerate}
    \item \textbf{High-Precision Classes ($>95\%$ Top-1)}: Visually distinct sports with unique equipment or specialized environments, such as \textit{Archery}, \textit{Bowling}, \textit{Billiards}, \textit{Skiing}, and \textit{Surfing}, consistently achieve near-perfect classification ($97.2\%$--$100\%$).
    \item \textbf{Semantic Ambiguity Clusters}: The primary sources of classification confusion occur in pairs sharing nearly identical athletic posture transitions:
    \begin{itemize}
        \item \textit{BaseballPitch} vs. \textit{CricketBowling}: Both feature an overhand throwing motion on grass/dirt pitches. However, VideoMAE successfully resolves this ambiguity in $92\%$ of test instances by learning to attend to the ball delivery stride and arm extension radius.
        \item \textit{TennisSwing} vs. \textit{TableTennisShot}: Confusion stems from rapid forehand/backhand wrist strokes; resolution requires fine-grained spatiotemporal attention to table geometry and athlete torso stance.
    \end{itemize}
\end{enumerate}
